% MATH 100 mock midterm.
% Student version (default):
%   pdflatex -interaction=nonstopmode -halt-on-error -output-directory=output/pdf Lecture8/mock-mid-term.tex
% Instructor version: change \showsolutionsfalse to \showsolutionstrue.
\documentclass[11pt]{article}

\usepackage[a4paper,margin=18mm,headheight=15pt]{geometry}
\usepackage{amsmath,amssymb,mathtools}
\usepackage{array,booktabs}
\usepackage{enumitem}
\usepackage{xcolor}
\usepackage{fancyhdr}
\usepackage{iftex}
\ifPDFTeX
  \usepackage[UTF8,heading=false]{ctex}
\else
  \usepackage{fontspec}
  \usepackage{xeCJK}
  \setCJKmainfont{Microsoft YaHei}
\fi

\definecolor{solutionblue}{RGB}{0,76,170}
\newif\ifshowsolutions
\showsolutionsfalse

\newcommand{\R}{\mathbb{R}}
\newcommand{\vect}[1]{\mathbf{#1}}
\newcommand{\Span}{\operatorname{Span}}
\newcommand{\Col}{\operatorname{Col}}
\newcommand{\Nul}{\operatorname{Nul}}
\newcommand{\rank}{\operatorname{rank}}
\newcommand{\nullity}{\operatorname{nullity}}
\newcommand{\pts}[1]{\hfill\textbf{[#1 pts]}}
\newcommand{\answerline}{\ifshowsolutions\else\par\vspace{0.55in}\hrulefill\par\fi}
\newcommand{\workarea}[1]{\ifshowsolutions\else\vspace{#1}\fi}
\newcommand{\solution}[1]{%
  \ifshowsolutions
    \par\medskip
    {\color{solutionblue}\textbf{Solution.} #1}
    \par\medskip
  \fi}

\pagestyle{fancy}
\fancyhf{}
\lhead{MATH 100: Linear Algebra}
\rhead{Mock Midterm}
\cfoot{\thepage}

\setlist[enumerate,1]{leftmargin=*,itemsep=1.0em}
\setlist[enumerate,2]{leftmargin=2.0em,itemsep=0.25em,topsep=0.35em}

\begin{document}

\begin{center}
  {\LARGE\bfseries MATH 100: Linear Algebra}\par
  \vspace{0.25em}
  {\Large\bfseries Mock Midterm Examination}\par
  \vspace{0.6em}
  \ifshowsolutions
    {\large\color{solutionblue}\bfseries Instructor solution version}\par
  \fi
\end{center}

\noindent
\begin{tabular}{@{}p{.47\linewidth}p{.47\linewidth}@{}}
Name: \hrulefill & Student ID: \hrulefill \\
\end{tabular}

\vspace{0.8em}
\noindent
\begin{tabular}{@{}lccc@{}}
\toprule
Section & Questions & Points per question & Total \\
\midrule
I. Multiple choice & 10 & 2 & 20 \\
II. Brief answers & 5 & 2 & 10 \\
III. Exercises & 7 & 10 & 70 \\
\midrule
Total & 22 & & 100 \\
\bottomrule
\end{tabular}

\vspace{0.8em}
\noindent\textbf{Instructions}
\begin{itemize}[leftmargin=1.5em,itemsep=0.25em,topsep=0.35em]
  \item This is a closed-book mock examination.
  \item Write one answer for each multiple-choice question.
  \item Give a short justification in Section II.
  \item Show the main steps in Section III. Unsupported answers may receive limited credit.
  \item Simplify exact answers unless the question states otherwise.
\end{itemize}

\section*{Section I. Multiple choice \hfill 20 points}

Choose the best answer. Each question is worth 2 points.

\begin{enumerate}
\item The reduced augmented matrix of a system is
\[
\left[\begin{array}{ccc|c}
1&0&2&3\\
0&1&-1&4\\
0&0&0&0
\end{array}\right].
\]
How many solutions does the system have?
\begin{enumerate}[label=\Alph*.]
  \item No solution
  \item Exactly one solution
  \item Infinitely many solutions
  \item The information is insufficient
\end{enumerate}
\solution{\textbf{C.} The third variable is free and the system is consistent.}

\item Which set is guaranteed to be linearly dependent?
\begin{enumerate}[label=\Alph*.]
  \item Two nonzero vectors in $\R^3$
  \item Three vectors in $\R^4$
  \item Four vectors in $\R^3$
  \item Two vectors that are not scalar multiples of one another
\end{enumerate}
\solution{\textbf{C.} Four vectors in $\R^3$ give more vectors than entries.}

\item If $A$ is $3\times2$ and $B$ is $2\times4$, what is the size of $AB$?
\begin{enumerate}[label=\Alph*.]
  \item $2\times2$
  \item $2\times4$
  \item $3\times2$
  \item $3\times4$
\end{enumerate}
\solution{\textbf{D.} The inner dimensions match and the outer dimensions are $3\times4$.}

\item If $A$ and $B$ are invertible matrices of the same size, which identity is correct?
\begin{enumerate}[label=\Alph*.]
  \item $(AB)^{-1}=A^{-1}B^{-1}$
  \item $(AB)^{-1}=B^{-1}A^{-1}$
  \item $(AB)^{-1}=AB^{-1}$
  \item $(AB)^{-1}=BA^{-1}$
\end{enumerate}
\solution{\textbf{B.} The inverse reverses the order of the factors.}

\item Let $T(\vect x)=A\vect x$, where $A$ is an $m\times n$ matrix. The transformation maps onto $\R^m$ exactly when
\begin{enumerate}[label=\Alph*.]
  \item $A$ has a pivot in every row.
  \item $A$ has a pivot in every column.
  \item $A\vect x=\vect0$ has a free variable.
  \item the columns of $A$ are linearly independent.
\end{enumerate}
\solution{\textbf{A.} A pivot in every row makes $A\vect x=\vect b$ consistent for every $\vect b\in\R^m$.}

\item Suppose $\det A=2$. Matrix $B$ is obtained from $A$ by multiplying one row by $-3$ and then interchanging two rows. What is $\det B$?
\begin{enumerate}[label=\Alph*.]
  \item $-6$
  \item $6$
  \item $-12$
  \item $12$
\end{enumerate}
\solution{\textbf{B.} The scale contributes $-3$ and the interchange contributes $-1$, so $\det B=(-3)(-1)(2)=6$.}

\item Which condition guarantees that an $n\times n$ matrix $A$ is singular?
\begin{enumerate}[label=\Alph*.]
  \item The columns of $A$ span $\R^n$.
  \item $\det A\ne0$.
  \item $A\vect x=\vect0$ has a nontrivial solution.
  \item $A$ is row equivalent to $I_n$.
\end{enumerate}
\solution{\textbf{C.} A nontrivial null-space vector means the columns are dependent and $A$ is singular.}

\item Let $A=[\vect a_1\ \vect a_2\ \vect a_3\ \vect a_4]$. Row reduction shows that the pivot columns are columns 1 and 3. Which set is a basis for $\Col A$?
\begin{enumerate}[label=\Alph*.]
  \item $\{\vect a_1,\vect a_2\}$
  \item $\{\vect a_1,\vect a_3\}$
  \item The first and third columns of the reduced matrix
  \item $\{\vect a_2,\vect a_4\}$
\end{enumerate}
\solution{\textbf{B.} Use the corresponding pivot columns from the original matrix.}

\item A $4\times7$ matrix has rank $3$. What is its nullity?
\begin{enumerate}[label=\Alph*.]
  \item $1$
  \item $3$
  \item $4$
  \item $7$
\end{enumerate}
\solution{\textbf{C.} Rank-nullity gives $3+\nullity A=7$.}

\item Which transformation from $\R^2$ to $\R^2$ is linear?
\begin{enumerate}[label=\Alph*.]
  \item $T(x,y)=(x+1,y)$
  \item $T(x,y)=(x^2,y)$
  \item $T(x,y)=(x+2y,3x-y)$
  \item $T(x,y)=(xy,y)$
\end{enumerate}
\solution{\textbf{C.} It is multiplication by $\begin{bmatrix}1&2\\3&-1\end{bmatrix}$.}
\end{enumerate}

\newpage
\section*{Section II. Brief answers \hfill 10 points}

Give a short answer and enough justification to support it. Each question is worth 2 points.

\begin{enumerate}
\item For which value of $h$ is the system represented below consistent?
\[
\left[\begin{array}{ccc|c}
1&0&-2&1\\
0&1&3&4\\
0&0&0&h-5
\end{array}\right].
\]
\answerline
\solution{The system is consistent exactly when $h-5=0$, so $\boxed{h=5}$.}

\item Determine the number of pivot variables and free variables in the echelon matrix
\[
\begin{bmatrix}
1&2&0&-1&0&3\\
0&0&1&4&0&-2\\
0&0&0&0&1&5\\
0&0&0&0&0&0
\end{bmatrix}.
\]
\answerline
\solution{Pivot columns are $1,3,5$. There are $\boxed{3}$ pivot variables and $\boxed{3}$ free variables, in columns $2,4,6$.}

\item Determine whether
\[
C=\begin{bmatrix}7&-2\\-3&1\end{bmatrix}
\]
is the inverse of
\[
A=\begin{bmatrix}1&2\\3&7\end{bmatrix}.
\]
\answerline
\solution{Yes. Direct multiplication gives $AC=CA=I_2$.}

\item Let
\[
\vect v_1=\begin{bmatrix}1\\0\\1\end{bmatrix},\quad
\vect v_2=\begin{bmatrix}0\\1\\1\end{bmatrix},\quad
\vect v_3=\begin{bmatrix}1\\1\\0\end{bmatrix},\quad
\vect b=\begin{bmatrix}2\\3\\1\end{bmatrix}.
\]
Is $\vect b\in\Span\{\vect v_1,\vect v_2,\vect v_3\}$?
\answerline
\solution{Yes. $\vect b=0\vect v_1+\vect v_2+2\vect v_3$.}

\item Suppose $A$ and $C$ are invertible square matrices and
\[
ABC=P.
\]
Solve for $B$ without changing the order of the factors.
\answerline
\solution{Multiply by $A^{-1}$ on the left and $C^{-1}$ on the right: $\boxed{B=A^{-1}PC^{-1}}$.}
\end{enumerate}

\newpage
\section*{Section III. Exercises \hfill 70 points}

Show the main steps. Each exercise is worth 10 points.

\begin{enumerate}
\item Solve the linear system. Write the final answer as an ordered triple. \pts{10}
\[
\begin{cases}
x+2y-z=3,\\
2x+5y+z=8,\\
-x-y+2z=-1.
\end{cases}
\]
\workarea{2.3in}
\solution{Row reduction, or substitution, gives $y+3z=2$ and $y+z=2$. Hence $z=0$, $y=2$, and $x=-1$. Thus $\boxed{(x,y,z)=(-1,2,0)}$.}

\item Compute $AB$. \pts{10}
\[
A=\begin{bmatrix}1&-2&3\\0&4&-1\end{bmatrix},
\qquad
B=\begin{bmatrix}2&1\\-1&3\\4&0\end{bmatrix}.
\]
\workarea{2.2in}
\solution{
\[
AB=
\begin{bmatrix}
1(2)+(-2)(-1)+3(4) & 1(1)+(-2)(3)+3(0)\\
0(2)+4(-1)+(-1)(4) & 0(1)+4(3)+(-1)(0)
\end{bmatrix}
=\boxed{\begin{bmatrix}16&-5\\-8&12\end{bmatrix}}.
\]}

\newpage
\item Compute the determinant efficiently. State the row or column you use for the first expansion. \pts{10}
\[
D=\begin{bmatrix}
2&0&0&0\\
1&3&4&0\\
0&2&-1&5\\
0&0&0&4
\end{bmatrix}.
\]
\workarea{2.2in}
\solution{Expand first along row 1, then along the last column of the remaining matrix:
\[
\det D=2\cdot4\det\begin{bmatrix}3&4\\2&-1\end{bmatrix}
=8(-3-8)=\boxed{-88}.
\]}

\item Find $A^{-1}$ by row reducing $[A\mid I]$. \pts{10}
\[
A=\begin{bmatrix}
1&1&0\\
0&1&1\\
1&0&1
\end{bmatrix}.
\]
\workarea{2.5in}
\solution{Row reduction gives
\[
\boxed{A^{-1}=\frac12
\begin{bmatrix}
1&-1&1\\
1&1&-1\\
-1&1&1
\end{bmatrix}}.
\]
Multiplication verifies $AA^{-1}=I_3$.}

\newpage
\item Find an LU factorization $A=LU$, where $L$ is unit lower triangular. \pts{10}
\[
A=\begin{bmatrix}
2&4&-2\\
4&9&-3\\
-2&-3&7
\end{bmatrix}.
\]
\workarea{2.7in}
\solution{Use multipliers $m_{21}=2$, $m_{31}=-1$, and $m_{32}=1$. Then
\[
\boxed{
L=\begin{bmatrix}1&0&0\\2&1&0\\-1&1&1\end{bmatrix},
\qquad
U=\begin{bmatrix}2&4&-2\\0&1&1\\0&0&4\end{bmatrix}}
\]
and $LU=A$.}

\item Matrix $B$ is obtained from a square matrix $A$ using the following operations in order:
\begin{enumerate}[label=(\roman*)]
  \item interchange rows 1 and 3;
  \item multiply row 2 by $3$;
  \item replace row 1 by row 1 minus twice row 2.
\end{enumerate}
If $\det B=36$, find $\det A$. \pts{10}
\workarea{2.0in}
\solution{The interchange multiplies the determinant by $-1$, the scaling multiplies it by $3$, and the row replacement does not change it. Thus
\[
36=(-1)(3)\det A,
\qquad
\boxed{\det A=-12}.
\]}

\newpage
\item Consider the parameterized system $A(k)\vect x=\vect b$, where
\[
A(k)=\begin{bmatrix}
1&1&0\\
0&k-1&1\\
0&0&k+2
\end{bmatrix},
\qquad
\vect b=\begin{bmatrix}2\\1\\3\end{bmatrix}.
\]
\begin{enumerate}[label=(\alph*)]
  \item Find all values of $k$ for which the system has a unique solution. \hfill[3 pts]
  \item When $k=1$, classify the system and write all solutions in parametric vector form. \hfill[4 pts]
  \item When $k=-2$, classify the system and justify your answer. \hfill[3 pts]
\end{enumerate}
\workarea{2.5in}
\solution{
Because $A(k)$ is triangular,
\[
\det A(k)=(k-1)(k+2).
\]
\begin{enumerate}[label=(\alph*)]
  \item The system has a unique solution when $\boxed{k\ne1,-2}$.
  \item At $k=1$, the equations reduce to $x_1+x_2=2$ and $x_3=1$; the last equation repeats $3x_3=3$. There is one free variable, so there are infinitely many solutions:
  \[
  \boxed{\vect x=\begin{bmatrix}2\\0\\1\end{bmatrix}
  +t\begin{bmatrix}-1\\1\\0\end{bmatrix},\quad t\in\R}.
  \]
  \item At $k=-2$, the last row becomes $0=3$. The system is inconsistent, so it has no solution.
\end{enumerate}}
\end{enumerate}

\ifshowsolutions
\newpage
\section*{Compact answer key}
\begin{description}[leftmargin=2.5em]
  \item[MC:] C, C, D, B, A, B, C, B, C, C.
  \item[Brief:] $h=5$; 3 pivots and 3 free variables; yes; $\vect b=\vect v_2+2\vect v_3$; $B=A^{-1}PC^{-1}$.
  \item[Exercises:] $(-1,2,0)$; $\begin{bmatrix}16&-5\\-8&12\end{bmatrix}$; $-88$; inverse shown above; LU shown above; $-12$; unique for $k\ne1,-2$, infinite at $k=1$, none at $k=-2$.
\end{description}
\fi

\end{document}
